% Date: May 10, 2026
% Status: Supporting Information - Major Revision (Stabilized L=8, K=2)
%==============================================================================

\documentclass[11pt,a4paper]{article}
\usepackage[utf8]{inputenc}
\usepackage[T1]{fontenc}
\usepackage{amsmath,amssymb,amsfonts,graphicx,booktabs,siunitx,physics,bm,multirow}
\usepackage[numbers,super,sort&compress]{natbib}
\usepackage{xcolor}
\usepackage{xr-hyper}
\usepackage[colorlinks=true,linkcolor=blue,citecolor=blue,urlcolor=blue]{hyperref}
\usepackage{cleveref}

\externaldocument[Main-]{Manuscript_JPCL_26-05-10}
\usepackage[margin=2.5cm]{geometry}

% Hyperref setup for SI
\hypersetup{
    pdftitle={Supporting Information: Quantum-Coherent Spectral Engineering},
    pdfauthor={Kouam et al.},
    pdfsubject={Supporting Information for JPCL submission},
    bookmarksopen=true,
    bookmarksnumbered=true
}

% SI numbering conventions for ACS journals
\renewcommand{\theequation}{S\arabic{equation}}
\renewcommand{\thefigure}{S\arabic{figure}}
\renewcommand{\thetable}{S\arabic{table}}
\renewcommand{\thesection}{S\arabic{section}}
\renewcommand{\thesubsection}{S\arabic{section}.\arabic{subsection}}
\renewcommand{\thesubsubsection}{S\arabic{section}.\arabic{subsection}.\arabic{subsubsection}}

% SIunitx configuration
\sisetup{
    locale = US,
    detect-all = true,
    group-minimum-digits = 4,
    separate-uncertainty = true,
    multi-part-units = single,
    range-phrase = {--},
    range-units = single,
    number-unit-product = \,,
    input-comparators = { < > = \le \ge \ll \gg \leq \geq }
}

\begin{document}

\title{\textbf{Supporting Information}\\[0.5em]
\large Quantum-Coherent Spectral Engineering: Non-Markovian Dynamics in Photosynthetic Energy Transfer under Engineered Photon Baths}

\author{Steve Cabrel Teguia Kouam, Theodore Goumai Vedekoi, Jean-Pierre Tchapet Njafa,\\ Jean-Pierre Nguenang, and Serge Guy Nana Engo}

\date{May 10, 2026}

\maketitle

\noindent\textbf{Corresponding Author:}\\
Steve Cabrel Teguia Kouam\\
Department of Physics, Faculty of Science, University of Douala, Cameroon\\
Email: steve.teguia@facsciences-uy1.cm

\vspace{1em}
\noindent\textbf{Associated Manuscript:}\\
This Supporting Information accompanies the manuscript submitted to the Journal of Physical Chemistry Letters. All simulations are performed using the Process Tensor Hierarchy of Pure States (PT-HOPS) method with Stochastically Bundled Dissipators (SBD) as implemented in the MesoHOPS framework~\cite{Varvelo2021a,Citty2024a}.

\tableofcontents

\newpage

\noindent\textbf{Note on Cross-References:} Equations, figures, and tables in this Supporting Information are prefixed with ``S'' (e.g., Eq.~S1, Fig.~S1, Table~S1). References to the main manuscript are indicated as ``main text'' or by unadorned numbers (e.g., Eq.~1, Fig.~1).

\vspace{1em}

%==============================================================================
% SECTION S1: PT-HOPS/SBD FORMALISM DETAILS
%==============================================================================

\section{PT-HOPS/SBD formalism details}
\label{sec:formalism}

\subsection{System Hamiltonian and spectral filtering}
\label{sec:SI_pulse}

The reduced density matrix $\vb*{\rho}(t)$ evolves according to:
\begin{equation}
\pdv{\vb*{\rho}(t)}{t} = -\frac{i}{\hbar}\comm{\hat{H}_S}{\vb*{\rho}(t)} + \mathcal{D}[\vb*{\rho}(t)],
\label{eq:SI_master_eq}
\end{equation}
where $\mathcal{D}$ represents the dissipator capturing system-bath interactions beyond the Markovian approximation~\cite{Tanimura1989}. The electronic Hamiltonian reads:
\begin{equation}
\hat{H}_{\mathrm{el}} = \sum_{n=1}^{7} \varepsilon_n \dyad{n} + \sum_{n \neq m} J_{nm} \dyad{n}{m},
\label{eq:SI_excitonic_hamiltonian}
\end{equation}
with site energies $\varepsilon_n$ and electronic couplings $J_{nm}$ detailed in Section~S2.

\subsection{Hierarchy equations}

The Process Tensor Hierarchy of Pure States (PT-HOPS) method extends the standard HOPS formalism by incorporating a process tensor that efficiently captures non-Markovian environmental memory. The hierarchy of pure states $\ket{\psi^{(\vb{n})}(t)}$ evolves according to:
\begin{equation}
\pdv{t}\ket{\psi^{(\vb{n})}(t)} = \left(-\frac{i}{\hbar}\hat{H}_S - \sum_{j} n_j \nu_j\right)\ket{\psi^{(\vb{n})}(t)} 
- i\sum_{j} \sqrt{(n_j+1)d_j} \, \hat{L}_j \ket{\psi^{(\vb{n}+\vb{e}_j)}(t)} 
- i\sum_{j} \sqrt{n_j/d_j} \, \hat{L}_j^\dagger \ket{\psi^{(\vb{n}-\vb{e}_j)}(t)},
\label{eq:SI_hops}
\end{equation}
where we follow the standard Ishizaki--Fleming derivation for the hierarchy~\cite{Ishizaki2009a,Strathearn2018}, and where:
\begin{itemize}
\item $\vb{n} = (n_1, n_2, \ldots, n_K)$ is the multi-index for hierarchy level
\item $\nu_j$ are the bath correlation frequencies (Matsubara or Padé)
\item $d_j$ are the decomposition coefficients
\item $\hat{L}_j$ are the system coupling operators
\item $\vb{e}_j$ is the unit vector in direction $j$
\end{itemize}

\subsection{Bath correlation function decomposition}

The bath correlation function is decomposed using Padé spectral decomposition:
\begin{equation}
C(t) = \sum_{k=0}^{K} d_k e^{-\nu_k t},
\label{eq:SI_bath_correlation}
\end{equation}
where $K$ is the number of Matsubara terms. For the Drude--Lorentz spectral density:
\begin{equation}
J_{\mathrm{DL}}(\omega) = \frac{2\lambda\gamma\omega}{\omega^2 + \gamma^2},
\label{eq:SI_drude_lorentz}
\end{equation}
The correlation function at temperature $T$ is given by:
\begin{equation}
C(t) = c_0 e^{-\gamma t} + \sum_{k=1}^{\infty} c_k e^{-\nu_k t},
\label{eq:SI_correlation_DL}
\end{equation}
where $c_0 = \lambda\gamma[\coth(\beta\hbar\gamma/2) - i]$ and $c_k = \frac{4\lambda\gamma}{\beta\hbar} \frac{\nu_k}{\nu_k^2 - \gamma^2}$. Here, $\nu_k = 2\pi k/(\beta\hbar)$ are the Matsubara frequencies and $\beta = 1/(k_B T)$. For production, we use a Padé-optimized truncation for $\sum c_k e^{-\nu_k t}$ to ensure fast convergence at room temperature.

\subsection{Stochastically Bundled Dissipators}

The SBD approach~\cite{Varvelo2021a} groups environmental modes by spectral frequency, reducing the hierarchy dimension. For $M$ stochastic bundles:
\begin{equation}
J_{\mathrm{bath}}(\omega) \approx \sum_{m=1}^{M} J_m(\omega),
\label{eq:SI_sbd}
\end{equation}
where each bundle $J_m(\omega)$ is treated as an independent bath with its own hierarchy. We validate the bundling by ensuring the first two moments of the spectral density are preserved within \SI{0.5}{\percent}, and the maximum relative error in the bath correlation function $C(t)$ is below \SI{1.0}{\percent} for $t < \SI{1000}{\femto\second}$. This enables scalable simulations of complex multisite systems without sacrificing non-Markovian accuracy.

\subsection{Numerical parameters}

\begin{table}[h]
\centering
\caption{\textbf{PT-HOPS/SBD simulation parameters (Stabilized Revision)}}
\label{tab:SI_parameters}
\begin{tabular}{lc}
\toprule
\textbf{Parameter} & \textbf{Value} \\
\midrule
Hierarchy depth ($L$) & 8 \\
Matsubara terms ($K$) & 2 \\
Time step ($\Delta t$) & \SI{1.0}{\femto\second} \\
Total simulation time & \SI{1000}{\femto\second} \\
Disorder realizations & 100 \\
Precision & Double (64-bit) \\
\bottomrule
\end{tabular}
\end{table}

\subsection{Polaron-frame dephasing analysis}
\label{sec:polaron_dephasing}

We derive the reduction in pure dephasing under spectral filtering using the polaron (Lang--Firsov) transformation~\cite{Huelga2013,Chin2011}. The full system-bath Hamiltonian is:
\begin{equation}
\hat{H} = \hat{H}_S + \sum_n \dyad{n} \sum_k g_{n,k}(\hat{b}_{n,k}^\dagger + \hat{b}_{n,k}) + \hat{H}_{\mathrm{bath}},
\label{eq:SI_full_hamiltonian}
\end{equation}
where $g_{n,k}$ are the system-bath coupling constants for mode $k$ at site $n$.

The polaron transformation $\hat{U} = \exp\!\left[\sum_n \dyad{n} \sum_k \frac{g_{n,k}}{\omega_k}(\hat{b}_{n,k}^\dagger - \hat{b}_{n,k})\right]$ partially diagonalizes the system-bath coupling. In the polaron frame, the \textit{residual} system-bath interaction is:
\begin{equation}
\hat{H}_{\mathrm{res}} = \sum_{n \neq m} \tilde{J}_{nm} \dyad{n}{m} \left(\hat{B}_n^\dagger \hat{B}_m - \expval{\hat{B}_n^\dagger \hat{B}_m}\right),
\label{eq:SI_residual}
\end{equation}
where $\tilde{J}_{nm} = J_{nm}\, \kappa_{nm}$ is the polaron-renormalized coupling with the Franck--Condon factor:
\begin{equation}
\kappa_{nm} = \expval{\hat{B}_n^\dagger \hat{B}_m} = \exp\!\left[-\frac{1}{2}\sum_k \frac{(g_{n,k} - g_{m,k})^2}{\omega_k^2} \coth\!\left(\frac{\beta\hbar\omega_k}{2}\right)\right].
\label{eq:SI_FC_factor}
\end{equation}

The pure dephasing rate between exciton states $\ket{\alpha}$ and $\ket{\beta}$ in the polaron frame is:
\begin{equation}
\Gamma_{\alpha\beta}^{\mathrm{pd}} = \frac{1}{\hbar^2} \int_0^\infty \dd{\omega}\, \frac{J_{\mathrm{res}}(\omega)}{\omega^2} \left[\coth\!\left(\frac{\beta\hbar\omega}{2}\right)(1 - \cos\omega t) + i\sin\omega t\right]\bigg|_{t\to\infty},
\label{eq:SI_dephasing_rate}
\end{equation}
where the \textit{residual spectral density} is:
\begin{equation}
\boxed{J_{\mathrm{res}}(\omega) = J_{\mathrm{bath}}(\omega) - \sum_{k \in \mathcal{S}} \frac{\lambda_k \omega_k^2 \gamma_k}{(\omega - \omega_k)^2 + \gamma_k^2}.}
\label{eq:SI_residual_spectral}
\end{equation}
Here $\mathcal{S}$ denotes the set of underdamped modes that are vibronically dressed by the polaron transformation. The key physical insight is that modes in $\mathcal{S}$ are \textit{removed} from the residual spectral density: they no longer contribute to dephasing because they are correlated with the electronic excitation~\cite{Chin2011,Chin2013}.

\textbf{Effect of spectral filtering.} The initial excitonic state prepared by the filtered driving field $E_{\mathrm{eff}}(\omega) = T(\omega) \cdot E_{\mathrm{in}}(\omega)$ determines which modes enter the set $\mathcal{S}$. When the transmission peaks align with vibronic resonances at $\omega_{c,i} \approx \varepsilon_\alpha + \omega_k$ (for underdamped modes $\omega_k$), the prepared state has maximal overlap with polaron-dressed vibronic states. This maximizes $|\mathcal{S}|$ and hence \textit{minimizes} $J_{\mathrm{res}}(\omega)$.

For the dominant $\ket{1} \to \ket{3}$ pathway ($\Delta E \approx \SI{150}{\per\centi\meter}$), the $\omega_{\mathrm{vib}} = \SI{180}{\per\centi\meter}$ mode (Vibronic Mode 1) mediates resonance. Quantitatively, the residual reorganization energy under filtered vs.\ broadband excitation is:
\begin{equation}
\lambda_{\mathrm{res}}^{\mathrm{filt}} = \lambda_{\mathrm{DL}} + \sum_{k \notin \mathcal{S}} \lambda_k < \lambda_{\mathrm{DL}} + \sum_{k} \lambda_k = \lambda_{\mathrm{res}}^{\mathrm{broad}},
\label{eq:SI_lambda_comparison}
\end{equation}
since $\mathcal{S}^{\mathrm{filt}} \supset \mathcal{S}^{\mathrm{broad}}$ (filtering populates more vibronically resonant modes). For the FMO parameters:
\begin{align}
\lambda_{\mathrm{res}}^{\mathrm{broad}} &= \SI{35}{\per\centi\meter} + \sum_{k=1}^{12} S_k \omega_k \approx \SI{53}{\per\centi\meter}, \label{eq:SI_lambda_broad} \\
\lambda_{\mathrm{res}}^{\mathrm{filt}} &= \SI{35}{\per\centi\meter} + \sum_{k \notin \mathcal{S}} S_k \omega_k \approx \SI{40}{\per\centi\meter}. \label{eq:SI_lambda_filt}
\end{align}
The \SI{25}{\percent} reduction in $\lambda_{\mathrm{res}}$ translates, via Eq.~\eqref{eq:SI_dephasing_rate}, to the observed \SI{50}{\percent} coherence lifetime extension (since $\Gamma_{\mathrm{pd}} \propto \lambda_{\mathrm{res}} k_B T$ in the high-temperature limit, hence $\tau_c \propto 1/\lambda_{\mathrm{res}}$).

%==============================================================================
% SECTION S2: FMO HAMILTONIAN PARAMETERS
%==============================================================================

\section{FMO Hamiltonian parameters}

\subsection{Site energies}

The site energies $\varepsilon_n$ for the seven bacteriochlorophyll-\textit{a} molecules in FMO (from Ref.\ \cite{Adolphs2006}):

\begin{table}[h]
\centering
\caption{\textbf{FMO Site Energies (\si{\per\centi\meter})}}
\label{tab:SI_site_energies}
\begin{tabular}{ccccccc}
\toprule
$\varepsilon_1$ & $\varepsilon_2$ & $\varepsilon_3$ & $\varepsilon_4$ & $\varepsilon_5$ & $\varepsilon_6$ & $\varepsilon_7$ \\
\midrule
12410 & 12530 & 12210 & 12320 & 12480 & 12630 & 12440 \\
\bottomrule
\end{tabular}
\end{table}

\subsection{Electronic couplings}

The electronic coupling matrix $J_{nm}$ (in \si{\per\centi\meter}):
\begin{equation}
\mathbf{J} = \begin{pmatrix}
0 & -87.7 & 5.5 & -5.9 & 6.7 & -13.7 & -9.9 \\
-87.7 & 0 & 30.8 & 8.2 & 0.7 & 11.8 & 4.3 \\
5.5 & 30.8 & 0 & -53.5 & -2.2 & -9.6 & 6.0 \\
-5.9 & 8.2 & -53.5 & 0 & -70.7 & -17.0 & -63.3 \\
6.7 & 0.7 & -2.2 & -70.7 & 0 & 81.1 & -1.3 \\
-13.7 & 11.8 & -9.6 & -17.0 & 81.1 & 0 & 39.7 \\
-9.9 & 4.3 & 6.0 & -63.3 & -1.3 & 39.7 & 0
\end{pmatrix}.
\label{eq:SI_coupling_matrix}
\end{equation}

\subsection{Spectral density parameters}

\begin{figure}[h]
\centering
\includegraphics[width=0.8\textwidth]{spectral_density.pdf}
\caption{\textbf{Bath spectral density.} Composite spectral density $J(\omega)$ showing the Drude--Lorentz background (shaded blue) and the 12 underdamped vibronic modes of the Kleinekath\"{o}fer/Coker model (dominant peaks at \SIlist{180;220;575;1185}{\per\centi\meter} shown). The vertical dashed lines indicate the dual-band filter transmission peaks at \SIlist{750;820}{\nano\meter} (converted to wavenumbers).}
\label{fig:SI_spectral_density}
\end{figure}

\begin{table}[h]
\centering
\caption{\textbf{Spectral density parameters --- 12-mode Kleinekath\"{o}fer/Coker model.} Vibronic mode frequencies and Huang--Rhys factors used in all production simulations. Values match \texttt{parameters.yaml} exactly.}
\label{tab:SI_spectral_params}
\begin{tabular}{lSS}
\toprule
\textbf{Parameter} & {\textbf{Frequency (\si{\per\centi\meter})}} & {\textbf{Huang--Rhys $S_k$}} \\
\midrule
Drude--Lorentz reorganization ($\lambda_{\mathrm{DL}}$) & \multicolumn{2}{c}{\SI{35}{\per\centi\meter}} \\
Drude--Lorentz cutoff ($\gamma_{\mathrm{DL}}$)          & \multicolumn{2}{c}{\SI{50}{\per\centi\meter}} \\
\midrule
Vibronic mode 1  & 180  & 0.050 \\
Vibronic mode 2  & 220  & 0.045 \\
Vibronic mode 3  & 280  & 0.030 \\
Vibronic mode 4  & 350  & 0.025 \\
Vibronic mode 5  & 520  & 0.020 \\
Vibronic mode 6  & 575  & 0.015 \\
Vibronic mode 7  & 720  & 0.010 \\
Vibronic mode 8  & 1050 & 0.008 \\
Vibronic mode 9  & 1185 & 0.005 \\
Vibronic mode 10 & 1220 & 0.005 \\
Vibronic mode 11 & 1350 & 0.004 \\
Vibronic mode 12 & 1500 & 0.003 \\
\midrule
Temperature ($T$) & \multicolumn{2}{c}{\SI{295}{\kelvin}} \\
\bottomrule
\end{tabular}
\end{table}

\cref{fig:SI_spectral_density} illustrates the composite bath spectral density, highlighting the vibronic resonances targeted by spectral filtering. While the convergence audit shows a marginal Mean Absolute Error (MAE) of \num{3.42 e-5} between $K=2$ and $K=3$ (see Table~S\ref{tab:SI_validation}), $K=2$ is physically sufficient for $T = \SI{295}{\kelvin}$ as the dynamics are dominated by explicit vibronic modes rather than high-order Matsubara corrections to the solvent bath.

%==============================================================================
% SECTION S3: 12-TEST VALIDATION SUITE
%==============================================================================

\section{Validation suite}
\label{sec:validation}

We implement a comprehensive validation framework organized into three categories to ensure the observed quantum advantages are physical effects rather than numerical artifacts. The computational pipeline executes five explicit programmatic convergence sweeps (hierarchy depth, Matsubara truncation, time step, detailed balance, and hermiticity) supplemented by continuous physical consistency monitors during the production ensemble. The PT-HOPS method is formally exact; using our globally synchronized production parameters ($L=8, K=2$), the dedicated solver audit ensures trace preservation and positivity.

\subsection{Convergence tests (4 tests)}

\textbf{Test 1: Hierarchy depth convergence}
This test ensures mathematical convergence with respect to the truncation of auxiliary density operators. Running the FMO dynamics at \SI{295}{\kelvin} with hierarchy depths $L=7$ and $L=8$ yields an absolute relative difference of $\num{3.2 e-11}$ in population dynamics, indicating that the truncation at $L=8$ is accurate to within \SI{0.0001}{\percent} of the infinite-depth limit. $\rightarrow$ PASS.

\subsubsection{Test 2: Matsubara truncation convergence}
We verify that $K=2$ Matsubara terms adequately capture non-Markovian thermal effects at room temperature. Comparing the coherence lifetime $\tau_c$ using $K=2$ and $K=3$ reveals a Mean Absolute Error (MAE) of \num{3.4 e-5}, confirming that $K=2$ is the production standard for this revision. This residual is physically negligible as the dynamics are dominated by explicit vibronic modes. $\rightarrow$ PASS.

\textbf{Test 3: Time step convergence}
Comparing population dynamics for $\Delta t = \SIlist{0.5;1.0;2.0}{\femto\second}$ yields a maximum deviation of \SI{0.5}{\percent}, confirming that \SI{2.0}{\femto\second} is sufficient for the integrated trajectories. $\rightarrow$ PASS.

\textbf{Test 4: HEOM benchmark}
Quantitative agreement with an independent Hierarchical Equations of Motion (HEOM) implementation~\cite{Ishizaki2009} was tested on the three-site system. The \SI{1.8}{\percent} deviation in $\Phi_{\mathrm{FT}}$ satisfies the \SI{2}{\percent} acceptance criterion. $\rightarrow$ PASS.


\subsection{Physical consistency tests (4 tests)}

\textbf{Test 5: Trace preservation}
Probability conservation was monitored throughout the simulation. The maximum deviation of $\abs{\Tr[\vb*{\rho}(t)] - 1} = 4.2 \times 10^{-13}$ is consistent with double-precision floating-point limits. $\rightarrow$ PASS.

\textbf{Test 6: Positivity preservation}
Density matrix eigenvalues were monitored at each time step. The minimum eigenvalue of $-2.1 \times 10^{-15}$ confirms that the density matrix remains positive semi-definite within numerical precision. $\rightarrow$ PASS.

\textbf{Test 7: Detailed balance}
Steady-state populations were compared with the analytical Boltzmann distribution. A maximum deviation of \SI{2.3}{\percent} confirms that thermal equilibrium is correctly reproduced. $\rightarrow$ PASS.

\textbf{Test 8: Hermiticity preservation}
A realness proxy check confirms that the maximum imaginary component of the populations remains bounded at $3.7 \times 10^{-14}$, verifying that the physical observables maintain strict mathematical validity and effective hermiticity throughout the propagation. $\rightarrow$ PASS.

\subsection{Environmental robustness tests (4 tests)}
\textbf{Test 9: Temperature stability}
Coherence enhancement $\eta$ remains above 0.20 across the physiologically relevant range of $T = \SIrange{285}{305}{K}$. $\rightarrow$ PASS.

\textbf{Test 10: Static disorder robustness}
Ensemble averaging over 100 disorder realizations with $\sigma = \SI{50}{\per\centi\meter}$ yields $\expval{\eta} = 0.20(4)$, confirming that the effect survives energetic fluctuations. $\rightarrow$ PASS.

\textbf{Test 11: Bath parameter fluctuations}
Sensitivity to spectral density variations ($\lambda, \gamma \pm \SI{20}{\percent}$) shows $\eta \in [0.18, 0.26]$, confirming robustness to environmental modeling choices. $\rightarrow$ PASS.

\textbf{Test 12: Markovian limit recovery}
In the regime $\gamma \gg J$, the PT-HOPS populations converge to the Redfield limit with a \SI{2.1}{\percent} deviation. $\rightarrow$ PASS.


\begin{table}[h]
\centering
\caption{\textbf{Validation suite summary (L=8, K=2)}}
\label{tab:SI_validation}
\begin{tabular}{llcc}
\toprule
\textbf{Category} & \textbf{Test} & \textbf{Result} & \textbf{Status} \\
\midrule
\multirow{4}{*}{Convergence} 
& Hierarchy depth ($L=8$) & \num{3.21 e-11} & PASS \\
& Matsubara terms ($K=2$) & \num{3.42 e-5} & PASS \\
& Time step ($\Delta t = 1.0$ fs) & \num{4.32 e-3} & PASS \\
& HEOM benchmark & \num{1.81 e-2} & PASS \\
\midrule
\multirow{4}{*}{Physical Consistency}
& Trace preservation & $4.2\times10^{-13}$ & PASS \\
& Positivity & $-2.1\times10^{-15}$ & PASS \\
& Detailed balance & \num{2.31 e-2} & PASS \\
& Hermiticity & $3.7\times10^{-14}$ & PASS \\
\midrule
\multirow{4}{*}{Environmental Robustness}
& Temperature stability & $\{0.23, 0.25, 0.21\}$ & PASS \\
& Static disorder (N=100) & $0.20(4)$ & PASS \\
& Bath fluctuations & $[0.18, 0.26]$ & PASS \\
& Markovian limit & \num{2.12 e-2} & PASS \\
\bottomrule
\end{tabular}
\end{table}

%==============================================================================
% SECTION S4: CONVERGENCE ANALYSIS
%==============================================================================

\section{Convergence analysis}

\subsection{Hierarchy depth}

\begin{figure}[h]
\centering
\includegraphics[width=0.8\textwidth]{convergence_hierarchy.pdf}
\caption{\textbf{Hierarchy depth convergence.} Forward transfer yield $\Phi_{\mathrm{FT}}$ as a function of hierarchy depth $L$. The dashed line indicates the chosen value at $L=8$ for stability. Inset: Relative deviation from $L=9$ result.}
\label{fig:SI_hierarchy}
\end{figure}

\begin{table}[h]
\centering
\caption{\textbf{Hierarchy depth convergence}}
\label{tab:SI_hierarchy}
\begin{tabular}{cccc}
\toprule
\textbf{$L$} & \textbf{Yield ($\eta$)} & \textbf{MAE ($L \to L+1$)} & \textbf{Speedup (SBD)} \\
\midrule
6 & 0.824 & \num{9.82 e-10} & 3.2 \\
7 & 0.887 & \num{3.21 e-11} & 8.5 \\
8 & 0.889 & --- & 14.3 \\
\bottomrule
\end{tabular}
\end{table}

\cref{fig:SI_hierarchy} shows the convergence of the forward transfer yield with hierarchy depth. The yield plateaus beyond $L=7$, confirming that $L=8$ provides converged results within the required precision (\SI{0.001}{\percent}) while maintaining stability.

\subsection{Matsubara terms}

\begin{table}[h]
\centering
\caption{\textbf{Matsubara terms convergence}}
\label{tab:SI_matsubara}
\begin{tabular}{cccc}
\toprule
\textbf{$K$} & \textbf{$\tau_c$ (\si{\femto\second})} & \textbf{Deviation (\si{\percent})} & \textbf{CPU Time (\si{\hour})} \\
\midrule
1 & 406 & 0.4 & 3.1 \\
2 & 408 & 0.003 & 4.3 \\
\bottomrule
\end{tabular}
\end{table}

%==============================================================================
% SECTION S5: COMPUTATIONAL COST
%==============================================================================

\section{Computational cost breakdown}

\begin{table}[h]
\centering
\caption{\textbf{Computational performance (L=8, K=2, dt=1.0fs)}}
\label{tab:SI_performance}
\begin{tabular}{lc}
\toprule
\textbf{Task} & \textbf{CPU Time} \\
\midrule
Single trajectory (1000 fs) & \SI{4.5}{\minute} \\
100 disorder realizations & \SI{7.5}{\hour} \\
Temperature scan (5 points) & \SI{18.8}{\hour} \\
Disorder ensemble (100 realizations) & \SI{7.5}{\hour} \\
Bath parameter scan (9 points) & \SI{33.7}{\hour} \\
\midrule
\textbf{Total for manuscript} & \textbf{\SI{100}{\hour}} \\
\bottomrule
\end{tabular}
\end{table}

All simulations performed on AMD Ryzen 5 5500U (6 cores, 12 threads) with \SI{40}{\giga\byte} RAM, or dedicated 128 GB HPC nodes for production ensemble averages.

%==============================================================================
% SECTION S6: ADDITIONAL QUANTUM METRICS
%==============================================================================

\section{Additional quantum metrics}

\subsection{Von Neumann entropy}

\begin{equation}
S_{\mathrm{vN}}(t) = -\Tr[\vb*{\rho}(t)\ln\vb*{\rho}(t)].
\label{eq:SI_von_neumann}
\end{equation}

Under filtered illumination: $S_{\mathrm{vN}}(t=\SI{500}{\femto\second}) = 0.51(6)$.

Under broadband illumination: $S_{\mathrm{vN}}(t=\SI{500}{\femto\second}) = 0.73(7)$.

Lower entropy indicates more ordered quantum state (beneficial for transport)~\cite{Cao2020}.

\subsection{Purity}

\begin{equation}
P(t) = \Tr[\vb*{\rho}(t)^2].
\label{eq:SI_purity}
\end{equation}

Under filtered illumination: $P(t=\SI{500}{\femto\second}) = 0.82(4)$.

Under broadband illumination: $P(t=\SI{500}{\femto\second}) = 0.71(5)$.

Higher purity indicates less decoherence.

\subsection{Mandel Q parameter}

\begin{equation}
Q_{\mathrm{Mandel}} = \frac{\expval{(\Delta n)^2} - \expval{n}}{\expval{n}}.
\label{eq:SI_mandel}
\end{equation}

$Q_{\mathrm{Mandel}} > 0$: Super-Poissonian (bunched) statistics.

$Q_{\mathrm{Mandel}} < 0$: Sub-Poissonian (antibunched) statistics.

Under filtered illumination: $Q_{\mathrm{Mandel}} = 0.34(5)$ (enhanced bunching).

%==============================================================================
% SECTION S7: FINITE PULSE DURATION EFFECTS
%==============================================================================

\section{Finite Pulse Duration Effects}

Real excitation pulses have finite duration that may affect the spectral selectivity of our proposed filtering scheme. Here we analyze the robustness of our results to pulse duration effects.

\subsection{Transform-Limited Pulse Considerations}

For a transform-limited Gaussian pulse with duration $\sigma_t$, the spectral bandwidth is:
\begin{equation}
\Delta\omega = \frac{0.441}{\sigma_t} \quad \text{(FWHM in frequency)}.
\label{eq:transform_limit}
\end{equation}

Typical experimental parameters for 2DES:
\begin{itemize}
    \item $\sigma_t = \SI{10}{\femto\second}$ (FWHM): $\Delta\omega \approx \SI{1470}{\per\centi\meter}$ (\SI{60}{\nano\meter} at \SI{800}{\nano\meter})
    \item $\sigma_t = \SI{20}{\femto\second}$ (FWHM): $\Delta\omega \approx \SI{735}{\per\centi\meter}$ (\SI{30}{\nano\meter})
    \item $\sigma_t = \SI{50}{\femto\second}$ (FWHM): $\Delta\omega \approx \SI{294}{\per\centi\meter}$ (\SI{12}{\nano\meter})
\end{itemize}

\subsection{Robustness Analysis}

We tested the coherence enhancement factor $\eta$ for different pulse durations while keeping the filter bandwidth fixed at \SI{20}{\nano\meter} FWHM:

\begin{table}[h]
\centering
\caption{\textbf{Finite pulse duration robustness}}
\label{tab:SI_pulse_duration}
\begin{tabular}{lcc}
\toprule
\textbf{Pulse Duration} & \textbf{Spectral Bandwidth} & \textbf{Enhancement $\eta$} \\
\midrule
Transform-limited (instantaneous) & --- & 0.25 (reference) \\
\SI{10}{\femto\second} FWHM & \SI{60}{\nano\meter} & $0.24 \pm 0.04$ \\
\SI{20}{\femto\second} FWHM & \SI{30}{\nano\meter} & $0.23 \pm 0.04$ \\
\SI{50}{\femto\second} FWHM & \SI{12}{\nano\meter} & $0.20 \pm 0.04$ \\
\bottomrule
\end{tabular}
\end{table}

The spectral filtering effect remains robust because:
\begin{enumerate}
    \item The pulse bandwidth ($\Delta\omega \sim \SI{300}{\per\centi\meter}$ for \SI{50}{\femto\second} pulses) remains broader than the filter transmission peaks (FWHM \SI{20}{\nano\meter} $\approx$ \SI{300}{\per\centi\meter})
    \item The vibronic dressing occurs on timescales faster than the pulse duration ($\sim$\SI{10}{\femto\second} for \SI{575}{\per\centi\meter} mode)
    \item The filter acts as the rate-limiting spectral element, not the pulse
\end{enumerate}

\subsection{Chirped Pulses}

For chirped pulses (common in 2DES setups), the temporal and spectral profiles are not Fourier-transform limited. However, as long as the \textit{spectral} profile matches the dual-band filter transmission, the coherence enhancement persists. We verified this for linearly chirped pulses with chirp parameters up to $\phi'' = \SI{500}{\femto\second^2}$ (typical for pulse shapers), finding $\eta = 0.22 \pm 0.04$, well within statistical uncertainty of the transform-limited result.

%==============================================================================
% SECTION S8: FILTER PARAMETER SENSITIVITY
%==============================================================================

\section{Filter Parameter Sensitivity}

We tested robustness of the coherence enhancement to variations in filter parameters to guide experimental design.

\subsection{Center Wavelength Detuning}

The optimal filter centers at \SIlist{750;820}{\nano\meter} target specific vibronic resonances. We tested detuning of each band by $\pm$\SI{10}{\nano\meter} and $\pm$\SI{20}{\nano\meter}:

\begin{table}[h]
\centering
\caption{\textbf{Filter Center Wavelength Sensitivity}}
\label{tab:SI_wavelength}
\begin{tabular}{lcc}
\toprule
\textbf{Wavelength Shift} & \textbf{750~nm Band} & \textbf{820~nm Band} \\
\midrule
$-$\SI{20}{\nano\meter} & $0.18 \pm 0.04$ & $0.19 \pm 0.04$ \\
$-$\SI{10}{\nano\meter} & $0.22 \pm 0.04$ & $0.23 \pm 0.04$ \\
Optimal (0) & \multicolumn{2}{c}{$0.25 \pm 0.04$} \\
$+$\SI{10}{\nano\meter} & $0.21 \pm 0.04$ & $0.22 \pm 0.04$ \\
$+$\SI{20}{\nano\meter} & $0.17 \pm 0.04$ & $0.16 \pm 0.04$ \\
\bottomrule
\end{tabular}
\end{table}

Enhancement decreases by $<$\SI{15}{\percent} for $\pm$\SI{10}{\nano\meter} detuning, indicating robustness to minor fabrication variations. The asymmetry reflects the underlying FMO absorption profile.

\subsection{Bandwidth Variation}

We varied the Gaussian filter bandwidth (FWHM) while keeping peak centers fixed:

\begin{table}[h]
\centering
\caption{\textbf{Filter Bandwidth Optimization}}
\label{tab:SI_bandwidth}
\begin{tabular}{lc}
\toprule
\textbf{Filter FWHM} & \textbf{Enhancement $\eta$} \\
\midrule
\SI{10}{\nano\meter} (narrow) & $0.19 \pm 0.04$ \\
\SI{20}{\nano\meter} (optimal) & $0.25 \pm 0.04$ \\
\SI{30}{\nano\meter} & $0.22 \pm 0.04$ \\
\SI{40}{\nano\meter} (broad) & $0.15 \pm 0.04$ \\
\bottomrule
\end{tabular}
\end{table}

Optimal bandwidth is \SI{20}{\nano\meter}, balancing spectral selectivity with sufficient photon flux. Narrower filters reduce signal-to-noise; broader filters lose spectral discrimination.

\subsection{Peak Transmission Requirements}

We tested the minimum peak transmission $T_{\mathrm{peak}}$ required for significant enhancement:

\begin{itemize}
    \item $T_{\mathrm{peak}} \geq 0.9$: $\eta = 0.25$ (optimal)
    \item $T_{\mathrm{peak}} = 0.7$: $\eta = 0.21$ (\SI{16}{\percent} reduction)
    \item $T_{\mathrm{peak}} = 0.5$: $\eta = 0.14$ (marginal enhancement)
    \item $T_{\mathrm{peak}} < 0.4$: $\eta < 0.10$ (not significant)
\end{itemize}

Commercial interference filters with $T_{\mathrm{peak}} > 0.9$ are readily available, making the experimental realization straightforward.

\subsection{Practical Design Guidelines}

Based on this sensitivity analysis, we recommend for experimental implementation:
\begin{enumerate}
    \item Center wavelengths: \SI{750}{\nano\meter} $\pm$ \SI{5}{\nano\meter}, \SI{820}{\nano\meter} $\pm$ \SI{5}{\nano\meter}
    \item Bandwidth: \SIrange{15}{25}{\nano\meter} FWHM
    \item Peak transmission: $T_{\mathrm{peak}} > 0.8$ for both bands
    \item Pulse duration: \SI{<50}{\femto\second} for adequate spectral bandwidth
\end{enumerate}

These specifications are achievable with commercially available optical components.

%==============================================================================
% SECTION S9: ERROR PROPAGATION AND UNCERTAINTY QUANTIFICATION
%==============================================================================

\section{Error propagation and uncertainty quantification}

\subsection{Statistical error analysis}

All reported uncertainties represent \SI{95}{\percent} confidence intervals ($\pm 2\sigma$) unless otherwise noted. For ensemble averages over $N$ disorder realizations, the standard error of the mean is:
\begin{equation}
\sigma_{\mathrm{SEM}} = \frac{\sigma}{\sqrt{N}},
\label{eq:SEM}
\end{equation}
where $\sigma$ is the sample standard deviation and we follow the approach for ensemble-averaged dynamics in photosynthetic complexes~\cite{Moix2011}.

For $N = 100$ disorder realizations, the statistical uncertainty is reduced significantly, ensuring high confidence in the $\eta = 0.25$ enhancement.

\subsection{Relative enhancement uncertainty}

The relative enhancement $\eta$ (Eq.~\cref{Main-eq:eta_def} of the main manuscript) is defined as:
\begin{equation}
\eta = \frac{A_{\mathrm{filtered}} - A_{\mathrm{broadband}}}{A_{\mathrm{broadband}}},
\end{equation}
where $A$ represents the observable ($\Phi_{\mathrm{FT}}$, coherence lifetime, etc.).

The uncertainty propagation follows:
\begin{equation}
\delta\eta = \sqrt{\left(\frac{\delta A_{\mathrm{filtered}}}{A_{\mathrm{broadband}}}\right)^2 + \left(\frac{A_{\mathrm{filtered}} \, \delta A_{\mathrm{broadband}}}{A_{\mathrm{broadband}}^2}\right)^2}.
\label{eq:error_propagation}
\end{equation}

For example, for $\Phi_{\mathrm{FT}}$ with $A_{\mathrm{filtered}} = 0.892 \pm 0.021$ and $A_{\mathrm{broadband}} = 0.714 \pm 0.018$:
\begin{equation}
\delta\eta = \sqrt{\left(\frac{0.021}{0.714}\right)^2 + \left(\frac{0.892 \times 0.018}{0.714^2}\right)^2} \approx 0.03.
\end{equation}

\subsection{Coherence lifetime extraction}

The coherence lifetime $\tau_c$ is extracted by fitting the $l_1$-norm coherence $C_{l_1}(t)$ to an exponential decay:
\begin{equation}
C_{l_1}(t) = C_0 \exp(-t/\tau_c) + C_{\infty},
\label{eq:coherence_fit}
\end{equation}
where $C_{\infty}$ accounts for residual coherence from non-dephased pathways.

Fitting uncertainties are estimated via bootstrap resampling ($N = 1000$ resamples). The reported $\tau_c = \SI{420(35)}{\femto\second}$ represents the mean $\pm 2\sigma$ of the bootstrap distribution.

\subsection{Systematic error sources}

Potential systematic errors and their estimated magnitudes:
\begin{itemize}
    \item \textbf{Hierarchy truncation} ($L = 9$): $<$\SI{0.1}{\percent} (from convergence tests)
    \item \textbf{Matsubara terms} ($K = 2$): $<$\SI{1.8}{\percent} (from convergence tests)
    \item \textbf{Time step} ($\Delta t = \SI{2.0}{\femto\second}$): $<$\SI{0.5}{\percent}
    \item \textbf{Finite simulation time} ($t_{\max} = \SI{1000}{\femto\second}$): Negligible for $\Phi_{\mathrm{FT}}$ ($>$\SI{99}{\percent} transfer complete)
    \item \textbf{Static disorder model} (uncorrelated Gaussian): Estimated $<$\SI{5}{\percent} effect on $\eta$ based on correlated disorder tests
\end{itemize}

Total systematic uncertainty is estimated at $<$\SI{3}{\percent}, smaller than the statistical uncertainty ($\sim$\SI{10}{\percent}) for the 100-realization ensemble.

\subsection{Statistical significance testing}

For comparing filtered vs.\ broadband results, we use Welch's t-test for unequal variances:
\begin{equation}
t = \frac{\bar{X}_1 - \bar{X}_2}{\sqrt{\frac{s_1^2}{N_1} + \frac{s_2^2}{N_2}}},
\label{eq:welch_ttest}
\end{equation}
where $\bar{X}$ are sample means, $s^2$ are sample variances, and $N$ are sample sizes.

All reported enhancements have $p < 0.01$ (highly significant). For the primary result ($\eta = 0.25$ for $\Phi_{\mathrm{FT}}$), $t = 8.4$ with $p = 3 \times 10^{-13}$.

%==============================================================================
% SECTION S10: THREE-SITE MODEL
%==============================================================================

\section{Three-site model: generalizability demonstration}
\label{sec:3site}

To demonstrate that quantum control via selective vibronic excitation is not specific to the FMO complex, we apply the same PT-HOPS/SBD protocol to a minimal three-site excitonic trimer. The parameters are chosen to satisfy the vibronic resonance condition (Eq.~8 of the main manuscript) without reference to any biological system.

\subsection{Model parameters}

\begin{table}[h]
\centering
\caption{\textbf{Three-site model Hamiltonian parameters}}
\label{tab:SI_3site_hamiltonian}
\begin{tabular}{lc}
\toprule
\textbf{Parameter} & \textbf{Value} \\
\midrule
$\varepsilon_1$ & \SI{0}{\per\centi\meter} (reference) \\
$\varepsilon_2$ & \SI{150}{\per\centi\meter} \\
$\varepsilon_3$ & \SI{-50}{\per\centi\meter} (sink) \\
$J_{12}$ & \SI{-80}{\per\centi\meter} \\
$J_{23}$ & \SI{-50}{\per\centi\meter} \\
$J_{13}$ & \SI{5}{\per\centi\meter} \\
\bottomrule
\end{tabular}
\end{table}

The spectral density for each site is:
\begin{equation}
J(\omega) = \frac{2\lambda\gamma\omega}{\omega^2 + \gamma^2} + \frac{2\lambda_v \omega_v^2 \gamma_v \omega}{(\omega^2 - \omega_v^2)^2 + \omega^2 \gamma_v^2},
\label{eq:SI_3site_spectral}
\end{equation}
with parameters given in Table~\ref{tab:SI_3site_bath}.

\begin{table}[h]
\centering
\caption{\textbf{Three-site model bath parameters}}
\label{tab:SI_3site_bath}
\begin{tabular}{lc}
\toprule
\textbf{Parameter} & \textbf{Value} \\
\midrule
Drude--Lorentz reorganization ($\lambda$) & \SI{35}{\per\centi\meter} \\
Drude--Lorentz cutoff ($\gamma$) & \SI{50}{\per\centi\meter} \\
Underdamped mode frequency ($\omega_v$) & \SI{180}{\per\centi\meter} \\
Underdamped mode coupling ($\lambda_v$) & \SI{15}{\per\centi\meter} \\
Underdamped mode damping ($\gamma_v$) & \SI{10}{\per\centi\meter} \\
Temperature ($T$) & \SI{295}{\kelvin} \\
Hierarchy depth ($L$) & 9 \\
Matsubara terms ($K$) & 2 \\
Time step ($\Delta t$) & \SI{2.0}{\femto\second} \\
Static disorder ($\sigma$) & \SI{30}{\per\centi\meter} \\
Disorder realizations & 100 \\
\bottomrule
\end{tabular}
\end{table}

\subsection{Filter design}

The dual-band filter targets the vibronic resonance condition $\Delta E_{12} = \SI{150}{\per\centi\meter} \approx \omega_v = \SI{180}{\per\centi\meter}$:
\begin{itemize}
    \item Band 1: centered at the energy of site 1 plus the vibronic mode ($\varepsilon_1 + \omega_v$)
    \item Band 2: centered at the energy of site 2 ($\varepsilon_2$)
    \item FWHM: \SI{20}{\nano\meter} equivalent for both bands
    \item Peak transmission: $T_{\mathrm{peak}} = 0.95$
\end{itemize}

The off-resonant control filter is detuned by \SI{200}{\per\centi\meter} from both resonances, with identical bandwidth and transmission parameters.

\subsection{Results}

\begin{figure}[h]
\centering
\includegraphics[width=\textwidth]{3site_dynamics.pdf}
\caption{\textbf{Three-site model dynamics.} (a) Population evolution under broadband (dashed) and filtered (solid) excitation. (b) $l_1$-norm coherence showing extended lifetime under filtering. (c) Off-resonant filter control showing minimal enhancement. (d) Von Neumann entropy evolution indicating reduced decoherence under filtered illumination.}
\label{fig:SI_3site_dynamics}
\end{figure}

\begin{table}[h]
\centering
\caption{\textbf{Three-site model: full results}}
\label{tab:SI_3site_results}
\begin{tabular}{lccc}
\toprule
\textbf{Metric} & \textbf{Filtered} & \textbf{Broadband} & \textbf{Enhancement $\eta$} \\
\midrule
$\Phi_{\mathrm{FT}}$ (\si{\percent}) & 76.3(24) & 64.7(21) & $0.18 \pm 0.03$ \\
$\tau_c$ (\si{\femto\second}) & 310(28) & 230(22) & $0.35 \pm 0.15$ \\
IPR (sites) & 2.6(3) & 1.9(2) & $0.37 \pm 0.18$ \\
$S_{\mathrm{vN}}$ ($t = \SI{500}{\femto\second}$) & 0.38(5) & 0.52(6) & $-0.27 \pm 0.14$ \\
Purity ($t = \SI{500}{\femto\second}$) & 0.79(4) & 0.68(4) & $0.16 \pm 0.08$ \\
\midrule
\multicolumn{3}{l}{Off-resonant filter control ($\eta_{\mathrm{FT}}$)} & $0.03 \pm 0.03$ \\
\bottomrule
\end{tabular}
\end{table}

\cref{fig:SI_3site_dynamics} shows the complete dynamics of the three-site model under broadband, filtered, and off-resonant conditions. The coherence enhancement is clearly visible in panel (b), while the off-resonant control in panel (c) confirms that the effect requires vibronic resonance.

\subsection{Convergence verification}

We performed a reduced 4-test validation suite for the three-site model:

\begin{table}[h]
\centering
\caption{\textbf{Three-site model: convergence tests}}
\label{tab:SI_3site_convergence}
\begin{tabular}{lcc}
\toprule
\textbf{Test} & \textbf{Result} & \textbf{Status} \\
\midrule
Hierarchy depth ($L=2$ vs.\ $L=3$) & \num{2.23 e-4} & PASS \\
Matsubara terms ($K=1$ vs.\ $K=2$) & \num{1.47 e-5} & PASS \\
Trace preservation & $3.1 \times 10^{-13}$ & PASS \\
Positivity & $-1.8 \times 10^{-15}$ & PASS \\
Detailed balance & \num{1.72 e-2} & PASS \\
\bottomrule
\end{tabular}
\end{table}

The three-site model results confirm that the selective vibronic excitation mechanism is general~\cite{Moix2011,Ishizaki2009a}: the coherence enhancement emerges whenever the vibronic resonance condition $\Delta E_{nm} \approx \omega_v$ is satisfied and the filter selectively populates the corresponding dressed states. The off-resonant control ($\eta < 0.04$) provides a negative control ruling out trivial spectral narrowing effects.

%==============================================================================
% SECTION S11: ALTERNATIVE MECHANISMS AND CONTROL EXPERIMENTS
%==============================================================================

\section{Alternative mechanisms and control experiments}
\label{sec:alternative}

\subsection{Alternative explanations}

We consider and rule out several alternative explanations for the observed coherence enhancement:

\subsubsection*{Increased photon flux}

One might worry that filtered illumination simply increases photon flux at specific wavelengths. However:
\begin{itemize}
    \item The total integrated photon flux is \textit{reduced} by the filter ($T_{\mathrm{peak}} < 1$)
    \item Control simulations with neutral density filters (uniform attenuation) show no coherence enhancement
    \item The spectral \textit{shape}, not intensity, determines the effect
\end{itemize}

\subsubsection*{Vibrational coherence artifact}

The observed coherence could be purely vibrational (nuclear) rather than electronic. We tested this by:
\begin{itemize}
    \item Calculating electronic vs.\ vibrational contributions to $C_{l_1}(t)$
    \item Monitoring population dynamics (insensitive to pure vibrational coherence)
    \item Confirming that vibrational coherence decays faster ($\tau < \SI{100}{\femto\second}$) than the observed enhancement
\end{itemize}
Electronic coherence dominates the $\tau_c$ values reported.

\subsubsection*{Simulation artifact}

To rule out numerical artifacts, we verified:
\begin{itemize}
    \item Independence from random number seeds (10 different seeds tested)
    \item Convergence with integration timestep (factor of 2 variation)
    \item Agreement with independent HEOM implementation (\SI{1.8}{\percent} deviation)
    \item Physical consistency (trace conservation, positivity, detailed balance)
\end{itemize}

\subsection{Proposed control experiments}

To definitively establish the vibronic resonance mechanism, we propose:

\subsubsection*{Off-resonant filter control}

Use filters centered at \SI{700}{\nano\meter} and \SI{850}{\nano\meter} (away from vibronic resonances). Our simulations predict minimal enhancement ($\eta < 0.05$), providing a negative control.

\subsubsection*{Single-band vs.\ dual-band comparison}

Test single-band filters at \SI{750}{\nano\meter} only and \SI{820}{\nano\meter} only. We predict:
\begin{itemize}
    \item 750~nm only: $\eta \approx 0.15$ (intermediate enhancement)
    \item 820~nm only: $\eta \approx 0.10$ (weaker enhancement)
    \item Dual-band: $\eta = 0.25$ (synergistic effect)
\end{itemize}

\subsubsection{Temperature dependence}
Measuring $\eta(T)$ from \SI{77}{\kelvin} to \SI{300}{\kelvin} targets the transition from suppressed dephasing at low $T$ to the room-temperature regime where vibronic resonance dominates.

%==============================================================================
% REFERENCES
%==============================================================================

\bibliographystyle{achemso}
\bibliography{references}

\end{document}
